\documentclass[12pt]{article}
\usepackage{amsmath, amssymb}
\usepackage[margin=1in]{geometry}
\setlength{\parskip}{6pt}
\title{QUBO Formulation: Database Transaction Serializability Verification}
\date{}
\begin{document}
\maketitle

\subsection*{Database Transaction Serializability Verification}

\paragraph{Problem Statement.}
Given a set of $N$ candidate serial execution orders for a collection of database transactions, each candidate order $i$ is associated with a pre-computed violation cost $c_i \ge 0$ that quantifies the number of observed read/write dependency constraints violated by that order. The task is to select exactly one candidate serial order such that the total violation cost is minimized. A violation count of zero indicates that the execution is fully serializable under the supplied constraints.

\paragraph{Decision Variables.}
Let $N$ denote the number of candidate serial orders. For each candidate $i \in \{0, 1, \dots, N-1\}$, introduce a binary decision variable $x_i \in \{0,1\}$. The variable $x_i$ equals $1$ if and only if candidate order $i$ is selected, and $0$ otherwise. The selection is enforced by the requirement that exactly one variable takes the value $1$.

\paragraph{QUBO Derivation.}
\textbf{Step 1 --- Original Objective.} The problem is to minimize the total violation cost over the binary selection variables:
\begin{align}
\min_{x \in \{0,1\}^N} \sum_{i=0}^{N-1} c_i x_i.
\end{align}

\textbf{Step 2 --- Constraint Formulation.} The requirement that exactly one candidate order is selected is expressed as:
\begin{align}
\sum_{i=0}^{N-1} x_i = 1.
\end{align}

\textbf{Step 3 --- Penalty Term Expansion.} We convert the equality constraint into a quadratic penalty term by adding $\lambda \left( \sum_{i=0}^{N-1} x_i - 1 \right)^2$ to the objective, where $\lambda > 0$ is a penalty weight. Expanding the square yields:
\begin{align}
\lambda \left( \sum_{i=0}^{N-1} x_i - 1 \right)^2 &= \lambda \left( \sum_{i=0}^{N-1} x_i^2 + \sum_{i=0}^{N-1} \sum_{\substack{j=0 \\ j \neq i}}^{N-1} x_i x_j - 2 \sum_{i=0}^{N-1} x_i + 1 \right).
\end{align}
Using the idempotent property of binary variables, $x_i^2 = x_i$, the expression simplifies to:
\begin{align}
\lambda \left( \sum_{i=0}^{N-1} x_i + \sum_{i \neq j} x_i x_j - 2 \sum_{i=0}^{N-1} x_i + 1 \right) &= \lambda \left( \sum_{i \neq j} x_i x_j - \sum_{i=0}^{N-1} x_i + 1 \right).
\end{align}
Recognizing that $\sum_{i \neq j} x_i x_j = 2 \sum_{0 \le i < j \le N-1} x_i x_j$, the penalty contribution to the Hamiltonian becomes:
\begin{align}
\lambda \left( 2 \sum_{0 \le i < j \le N-1} x_i x_j - \sum_{i=0}^{N-1} x_i + 1 \right).
\end{align}

\textbf{Step 4 --- Combined Hamiltonian.} Adding the objective and the expanded penalty term gives the full QUBO Hamiltonian $H(x)$:
\begin{align}
H(x) = \sum_{i=0}^{N-1} c_i x_i + \lambda \left( 2 \sum_{0 \le i < j \le N-1} x_i x_j - \sum_{i=0}^{N-1} x_i + 1 \right).
\end{align}
Grouping the linear terms results in:
\begin{align}
H(x) = \sum_{i=0}^{N-1} (c_i - \lambda) x_i + 2\lambda \sum_{0 \le i < j \le N-1} x_i x_j + \lambda.
\end{align}

\textbf{Step 5 --- Q Matrix and Offset.} Comparing $H(x)$ with the standard QUBO form $H(x) = \sum_{i \le j} Q_{i,j} x_i x_j + c$, we read off the general closed-form expressions for the Q matrix entries and the constant offset:
\begin{align}
Q_{i,i} &= c_i - \lambda, & \forall i \in \{0, \dots, N-1\}, \\
Q_{i,j} &= \lambda, & \forall 0 \le i < j \le N-1, \\
c &= \lambda. &
\end{align}

\paragraph{Penalty Weight Justification.}
Let $C_{\max} = \max_{i} c_i$ denote the maximum violation cost among all candidates. If the constraint $\sum_{i=0}^{N-1} x_i = 1$ is violated, the squared term $\left( \sum_{i=0}^{N-1} x_i - 1 \right)^2$ evaluates to at least $1$ for any integer assignment, incurring a penalty of at least $\lambda$. The maximum possible reduction in the objective function by selecting a candidate is $C_{\max}$. Therefore, choosing $\lambda > C_{\max}$ ensures that the energy cost of violating the constraint strictly exceeds any potential gain in the objective function, guaranteeing that infeasible configurations are energetically unfavorable.

\paragraph{Correctness Argument.}
Minimizing $H(x)$ over $\{0,1\}^N$ recovers the optimal serial order because (i) for any feasible assignment where exactly one $x_i = 1$, the penalty term vanishes and $H(x)$ reduces exactly to the original objective $\sum_{i=0}^{N-1} c_i x_i$. (ii) For any infeasible assignment where $\sum_{i=0}^{N-1} x_i \neq 1$, the penalty term adds at least $\lambda$ to the energy. By enforcing $\lambda > C_{\max}$, the energy of every infeasible configuration is strictly greater than the energy of the best feasible configuration. Consequently, the global minimum of $H(x)$ must lie in the feasible subspace and correspond to the candidate order with the minimum violation cost.

\paragraph{Illustrative Example.}
Consider an instance with $N=3$ candidate orders and violation costs $c_0 = 5$, $c_1 = 2$, and $c_2 = 8$. We select $\lambda = 10$, which satisfies $\lambda > \max\{5, 2, 8\}$. The concrete objective is $5x_0 + 2x_1 + 8x_2$. The constraint is $x_0 + x_1 + x_2 = 1$. Substituting into the expanded penalty formula yields:
\begin{align}
10 \left( 2x_0x_1 + 2x_0x_2 + 2x_1x_2 - x_0 - x_1 - x_2 + 1 \right).
\end{align}
Combining with the objective gives the concrete Hamiltonian:
\begin{align}
H(x) = -5x_0 - 8x_1 - 2x_2 + 20x_0x_1 + 20x_0x_2 + 20x_1x_2 + 10.
\end{align}
The corresponding Q matrix entries and offset are:
\begin{align}
Q_{0,0} = -5, \quad Q_{1,1} = -8, \quad Q_{2,2} = -2, \quad Q_{0,1} = Q_{0,2} = Q_{1,2} = 10, \quad c = 10.
\end{align}
Evaluating $H(x)$ over the three feasible binary vectors:
\begin{align}
x = (1,0,0) \implies H = 5, \quad x = (0,1,0) \implies H = 2, \quad x = (0,0,1) \implies H = 8.
\end{align}
The global minimum occurs at $x^* = (0,1,0)$ with energy $H(x^*) = 2$, correctly identifying candidate order $1$ as the optimal selection with the minimum violation count.
\end{document}